\section{Arithmetic curves, inertia, sections, and faithfulness}
\label{sec:D}

Throughout this section a hyperbolic curve is a smooth geometrically connected
curve $U/K$ with smooth proper completion $X$, reduced complement $S=X\setminus U$,
and type $(g,n)=(g(X),\#S(\overline K))$ satisfying $2g-2+n>0$.
Unless a different coefficient set is explicitly specified, all fundamental groups
are full profinite \'etale groups.  We choose compatible geometric base points and
write
\[
 1\longrightarrow\Delta_U\longrightarrow\Pi_U
   \xrightarrow{p_U}G_K\longrightarrow1,
 \qquad \Delta_U=\pi_1(U_{\overline K}),\quad G_K=\operatorname{Gal}(\overline K/K).
\]
Changing a geometric base point changes maps by geometric inner conjugacy.
An inertia subgroup is the image of the geometric fundamental group of the
punctured strict henselian local scheme at a point of $S_{\overline K}$;
it is a conjugacy class of closed subgroups $I_s\simeq\widehat{\mathbb Z}(1)$.
We write $\operatorname{Sec}(U/K)$ for continuous sections of $p_U$, modulo
conjugation by $\Delta_U$, and $\rho_U:G_K\to\operatorname{Out}(\Delta_U)$ for
the outer action.  The topology on automorphism groups is the topology of
uniform agreement on finite characteristic quotients.

The distinction between an extension, a rational-point section of that extension,
and a section of a moduli-stack extension is essential in reading \S34.
We first settle the statements that do not depend on the three 2026 proof packs.

\subsection{What the arithmetic extension determines}

\statement{D.1}{Proved}{A centre-free extension is its outer-action pullback}
Let $\Delta$ be a topologically finitely generated centre-free profinite group
and let $1\to\Delta\to E\xrightarrow{p}G\to1$ be an extension of profinite groups.
If $\rho:G\to\operatorname{Out}(\Delta)$ is its associated outer action, then
\[
 E\longrightarrow
 \operatorname{Aut}(\Delta)\mathbin{\times}_{\operatorname{Out}(\Delta)}G,
 \qquad e\longmapsto(\operatorname{conj}_e|_\Delta,p(e))
\]
is an isomorphism of extensions.  Sections of $p$, modulo $\Delta$-conjugacy,
therefore correspond exactly to homomorphic lifts of $\rho$ to
$\operatorname{Aut}(\Delta)$, modulo $\operatorname{Inn}(\Delta)$-conjugacy.
This justifies the equivalence asserted in \S34, folio 245, scan 483.
\begin{proof}
The displayed map is a continuous homomorphism.  Its kernel lies in $\Delta$
and centralizes $\Delta$, hence is trivial.  Given $(a,g)$ in the fibre product,
choose $e\in E$ above $g$.  Equality of outer classes says
$a=\operatorname{conj}_d\operatorname{conj}_e$ for some $d\in\Delta$; thus
$(a,g)$ is the image of $de$.  A continuous bijection between compact Hausdorff
groups is an isomorphism.  Both assertions about sections follow from this
isomorphism.  Notice that $E\to\operatorname{Aut}(\Delta)$ by itself need not be
injective: it is the pair with $p$ that is always injective.
\end{proof}

\statement{D.2}{Refuted}{The literal open-normalizer criterion in scan 483}
It is false that the cuspidal inertia subgroups are the maximal closed subgroups
$L\subseteq\Delta_U$ for which $p_U(N_{\Pi_U}(L))$ is open in $G_K$.
It is also false that every nontrivial closed subgroup with this property is
contained in one inertia subgroup, as suggested in the margin of scan 483.
These failures already occur for $U=\mathbb P^1_{\mathbb Q}\setminus\{0,1,\infty\}$.
\begin{proof}
The subgroup $L=\Delta_U$ is normal in $\Pi_U$, so its normalizer is all of
$\Pi_U$ and its projected normalizer is $G_K$.  It is maximal among subgroups
with the stated property simply because it contains every subgroup under
consideration.  Here $\Delta_U\simeq\widehat F_2$ has quotient
$(\mathbb Z/2)^2$, whereas every inertia subgroup is procyclic.  Thus it cannot
be an inertia subgroup or be contained in one.  This is an example in the
arithmetic group of the curve, not merely an abstract profinite extension.

Even adding the word ``proper'' would not repair the marginal assertion.
The characteristic subgroup $[\Delta_U,\Delta_U]$ is proper, is normalized by
$\Pi_U$, and is not procyclic.  For the last assertion map the two free generators
onto $(12345)$ and $(123)$ in $A_5$.  They generate $A_5$: conjugates give
adjacent $3$-cycles, which generate all even permutations.  Since $A_5$ is
perfect, the image of $[\Delta_U,\Delta_U]$ is $A_5$, which is not cyclic.
The attached computation \texttt{code/d\_counterexamples.py} verifies these
finite calculations and records its output.

Finally, if the condition were instead imposed on subgroups of $\Pi_U$, every
rational-point section $s(G_K)$ would qualify because its normalizer contains
$s(G_K)$.  Such a subgroup projects isomorphically onto $G_K$ and is not
geometric inertia.  For instance $2\in U(\mathbb Q)$ gives such a section.
This explains why rational-point decomposition groups must not be silently
identified with cuspidal inertia.
\end{proof}
The scan actually says subgroups ``dans $\pi$''; no doubtful glyph supplies the
missing procyclic or cyclotomic hypotheses.  The alternatives
$\Gamma^+$ and $\Gamma$ in the projected-normalizer phrase make no difference:
$G_K$ is open in $G_{\mathbb Q}$.

\statement{D.3}{Known}{The cyclotomic recognition theorem for affine inertia}
Let $U/K$ be an affine hyperbolic curve over a number field.  A nontrivial
element $z\in\Delta_U$ belongs to a cuspidal inertia subgroup if and only if
there is a closed subgroup $N\subseteq\Pi_U$ such that $N$ normalizes
$\overline{\langle z\rangle}$, $p_U(N)$ is open in $G_K$, and
\[
 hzh^{-1}=z^{\chi_K(p_U(h))}\qquad(h\in N),
\]
where $\chi_K:G_K\to\widehat{\mathbb Z}^{\times}$ is the full cyclotomic
character.  In particular the action factors through $p_U(N)$; this condition
includes triviality of the conjugation action of $N\cap\Delta_U$.
The maximal procyclic subgroups having this property are precisely the
cuspidal inertia subgroups.  The collection is intrinsic to the arithmetic
extension together with its cyclotomic character.
\begin{proof}
The equivalence is \citet[Theorem (2.1.1)]{Nakamura1994}, applied to the full
class of finite groups.  The source assumes $n\geq1$, characteristic zero,
and a number field; all these hypotheses have been retained here.
For a cusp the branch-cycle description of its decomposition group gives
the displayed conjugation law.  Conversely Nakamura's weight argument identifies
$z$ as peripheral.  If a procyclic subgroup satisfying the condition is generated
topologically by $z$, it lies in the inertia subgroup supplied by the theorem.
The inertia subgroup itself satisfies the condition and is therefore the maximal
one.  For completeness, distinct maximal peripheral inertia subgroups intersect
trivially: if $1\ne x$ belongs to two of them, \citet[Proposition 3.6]{BoggiZalesskii2019BC}
gives $C_{\Delta_U}(x)=I$ for the first subgroup $I$, since a peripheral
generator is algebraically simple.  The second subgroup centralizes $x$, hence
is contained in $I$ and equals $I$ by maximality.  This also gives uniqueness
of the containing inertia subgroup for a nontrivial cyclotomic subgroup.
\end{proof}
We do not erase ``cyclotomic'' from this theorem.  The literal assertion refuted
in D.2 and the published assertion in D.3 have different hypotheses.  An
isomorphism of arithmetic extensions over $G_K$ preserves the latter hypotheses,
so it recovers the inertia conjugacy classes without their being supplied as
extra data.  The conclusion for arbitrary open homomorphisms follows more
cleanly from the Hom theorem below than by treating a normalizer calculation
as if it commuted with a noninjective homomorphism.

\statement{D.4}{Known}{The relative Hom theorem, including its exact ambiguity}
Let $K$ be sub-$p$-adic, meaning a subfield of a finitely generated extension
of $\mathbb Q_p$, and let $U,V$ be hyperbolic curves over $K$.  There is a
bijection
\[
 \operatorname{Hom}^{\mathrm{dom}}_K(U,V)
 \xrightarrow{\ \sim\ }
 \operatorname{Hom}^{\mathrm{open}}_{G_K}(\Pi_U,\Pi_V)/\operatorname{Inn}(\Delta_V).
\]
On the right the maps are continuous homomorphisms with open image whose
composite with $p_V$ is $p_U$.  In particular this applies over every number
field and every finite extension of $\mathbb Q_p$, including $p=2,3$.
A homomorphism in this set takes a cuspidal inertia subgroup either to the
trivial group or into a conjugate of a target cuspidal inertia subgroup.
\begin{proof}
The bijection is \citet[Theorem A; Theorem 16.5 and its following
Remark]{Mochizuki1999}.  We use the full profinite version.  The paper also has a
geometrically pro-$p$ version, whose arithmetic quotient is still the full
$G_K$; it is not a theorem about replacing the whole arithmetic group by its
maximal pro-$p$ quotient.

For clarity we verify the peripheral consequence, which is the assertion of
\S34, folios 245--245 verso, scans 483--484.  The associated nonconstant morphism
$f:U\to V$ extends uniquely to the smooth proper completions $\bar f:X\to Y$.
If a deleted point of $X$ maps to a point of $V$, its local inertia maps trivially
in $\pi_1(V)$.  If it maps to a deleted point of $Y$, choose local parameters
$t,u$ at the source and target.  The equality $\bar f^*u=t^e a$ with $a$ a local
unit shows that tame geometric inertia maps into the target inertia by
multiplication by the positive ramification index $e$.  Characteristic zero
makes all these horizontal inertia groups tame.  Thus ``respects loops''
allows a loop to be killed; it does not mean a permutation of all punctures.
\end{proof}
The relative Isom theorem follows formally: the two homomorphisms associated
to a group isomorphism and its inverse give morphisms whose composites are
identities by the injectivity of this bijection.  No separate compatibility
with a previously chosen loop system is required.

\statement{D.5}{Proved}{Admissible effective lifts are exactly the point sections}
Let $U/K$ be hyperbolic with $K$ sub-$p$-adic.  Call a lift of $\rho_U$
\emph{geometrically admissible} if it is obtained from a hyperbolic curve $V/K$,
a point $v\in V(K)$, and a $G_K$-equivariant outer isomorphism
$\Delta_V\simeq\Delta_U$ by transporting the action defined by $v$.
Then the admissible lifts modulo $\Delta_U$-conjugacy are in bijection with
$U(K)$.  This proves the deduction in \S34, scan 485, folio read
$246\mid245$, without assuming surjectivity of the section map.
\begin{proof}
By D.1 the given outer isomorphism produces an isomorphism
$\Pi_V\simeq\Pi_U$ over $G_K$, with the stated geometric ambiguity.
The Isom consequence of D.4 gives a unique $K$-isomorphism $f:V\to U$.
Functoriality carries the section at $v$ to the section at $f(v)$.
Conversely the point $u\in U(K)$ supplies the required admissible datum by
taking $V=U$, $v=u$, and the identity outer isomorphism.
Finally two points cannot give the same section class:
\citet[Theorem C = Theorem 19.1]{Mochizuki1999} proves injectivity of the
point-to-section map under precisely these field and hyperbolicity hypotheses.
Thus the two constructions give the asserted bijection.  Changing the chosen
representative of the outer isomorphism conjugates the lift by $\Delta_U$,
which is exactly the ambiguity already factored out.
\end{proof}
This proposition is about lifts that are already assumed admissible.  It does
not say that every lift is admissible.  For a proper curve of genus at least two
that additional assertion is the section conjecture; for an affine curve the
presence of cuspidal sections changes the assertion itself.

\statement{D.6}{Refuted}{All affine sections need not come from interior points}
For $U=\mathbb P^1_{\mathbb Q}\setminus\{0,1,\infty\}$ there is a continuous
section of $\Pi_U\to G_{\mathbb Q}$ which is not the section of any point of
$U(\mathbb Q)$.  Thus an unqualified affine version of ``every lift is the
lift of a rational point of $U$'' is false.
\begin{proof}
Take the rational tangential base point at $t=0$ with tangent vector $1$.
The associated fibre functor gives a section $s_0$.  More explicitly, embed the
local punctured disk in the Puiseux field
$\bigcup_m\overline{\mathbb Q}((t^{1/m}))$ and let $G_{\mathbb Q}$ act on
coefficients while fixing the compatible symbols $t^{1/m}$.  On the finite
covers of $U$ this gives the tangential section; changing paths conjugates it
by $\Delta_U$.

Under the morphism $t:U\to\mathbb G_m$, its Kummer class in
$H^1(G_{\mathbb Q},\mu_m)=\mathbb Q^\times/(\mathbb Q^\times)^m$ is zero for
every $m$, since each $t^{1/m}$ is fixed by the section.  If $s_0$ were the
section of $u\in U(\mathbb Q)$, the same classes would be the classes of $u$.
Hence $u$ would belong to $\bigcap_{m\geq1}(\mathbb Q^\times)^m$.
Every rational prime valuation of such a number is divisible by every $m$,
so all valuations vanish and $u=\pm1$.  Taking $m=2$ removes $-1$, giving
$u=1$, which is not in $U$.  This contradiction proves the assertion.
\end{proof}
The appropriate affine section conjecture includes rational cusps and their
cuspidal sections.  D.6 does not refute that conjecture.

\subsection{The section manuscript: an audited conditional chain}

\statement{D.7}{Open}{The 2026 proper $p$-adic section claim}
The assertion
\[
 X(k)\xrightarrow{\sim}\operatorname{Sec}(X/k)
 \quad(k/\mathbb Q_p\text{ finite},\ X/k\text{ proper},\ g(X)\geq2)
\]
for every prime $p$ is claimed in \citet[Theorem 1.1]{OpenAI019}.
The precise new input is its Theorem 1.2: for $F=k(X)$, a fixed separable
closure $F^{\mathrm{sep}}$, and the subfield $F^{\mathrm{et}}$ generated by all
connected finite \'etale covers of $X$ with all their $F$-embeddings,
\[
 F^{\mathrm{sep}}\subseteq\widehat{F^{\mathrm{et}}}
 \quad\hbox{inside }\widehat{F^{\mathrm{sep}}}
 \tag{D.7.1}
\]
for every fixed prolongation of every real-valued rank-one norm on $F$
extending the normalized $p$-adic norm.  The copies of the fields in this
inclusion are specified; an abstract embedding of a finite extension into
some unrelated completion would not suffice.
We retain the label Open for the unconditional new claim in this audit:
claimed, unrefereed, and not completely verified by us.  We have found no
counterexample.  The conditional implication D.8 is proved below.

\paragraph{What we checked in Sections 3--6 of the pack.}
We read these sections in full and followed their logical order, including
the distinction between horizontal ramification in characteristic zero and
wild reduction in characteristic $p$.
The patching construction in Proposition 3.2 requires a root on every
component above each overlap annulus.  With that hypothesis the pulled-back
branched cover splits on the overlap, so gluing trivial covers over the inner
disks produces a cover \'etale over the entire proper curve.  Proper GAGA and
invariance of finite \'etale covers under extension between algebraically
closed characteristic-zero fields are used after gluing.  Lemma 3.3 then
evaluates at the chosen completion-valued point.  We checked why this yields
a field embedding rather than merely an analytic map.  Lemma 3.4 uses cyclic
horizontal inertia over $\overline k((t))$: divisibility of the branch indices
by one integer is preserved under compositum and Galois closure, and an
upper cover \'etale over a lower cover introduces no further horizontal
ramification.

For Lemma 4.2 we checked the graph constructions separately.  A cycle in the
stable dual graph carries a tame cyclic cover of degree $\ell\ne p$, with
expansion factor $\ell$ on each lifted edge.  A separating edge is first
lifted to parallel edges using an unramified cover on the two sides; stability
ensures that each side has a cycle or a component of positive genus.
We checked that taking fibre products and then connected components retains
the required roots on all remaining sheets.  The external lifting result is
\citet[Theorem 7.7]{ABBR2015}; its tame-cover definition includes absence of
additional branch points when the marked puncture set is empty.
For Lemma 4.3 we opened \citet[Theorem A and Proposition 2.4(iv)]{MochizukiTsujimura2024}.
The latter gives a stable model of a finite \'etale cover dominating a prescribed
model.  Taking an irrational valuation level excludes vertices and smooth
disk tubes, so the inverse image of a sufficiently small rational-ended
collar lies in node annuli.  This is the specific model domination needed;
semistable reduction alone would not provide it.  The accessible primary
version is dated June 16, 2023; we did not silently identify it with the
March 30, 2024 revision cited by the pack.

In Section 5 the period lattice in a totally degenerate Jacobian is torsion-free.
Consequently local lifts of a multiplication-by-$m$ pullback can be adjusted
by periods and glued without a residual torsion ambiguity.  We checked that
Lemma 5.1 is intended to give a root of a torus character on the whole pullback,
not just one component.  Proposition 5.2 uses
$y^\ell=T(T-1)/(T-\alpha)$, with $\ell$ odd and $\ell\ne p$,
whose skeleton has two rational vertices joined by $\ell$ edges.
A cycle formed from two edges supplies a torus character with slope one on
the chosen edge.  If $u=cT(1+h)$ and $s=v(T)$ on an annulus of length $a$,
the Laurent coefficient inequalities give
$v(h)\geq\min(s,a-s)$.  The trimming condition
\[
 \min(s,a-s)>v_p(m)+1/(p-1)
\]
is sufficient for an analytic $m$th root of $1+h$, also when $p\mid m$.
We checked this numerical point; replacing it by a prime-to-$p$ root argument
would not prove Theorem 1.2.  We also checked the use of normality in Section 6:
a completion-valued embedding of a finite normal extension into the fixed
completion forces its specified image there to be the same normal subfield.

This is a substantial audit of the proof's mechanism, not a complete independent
verification of all analytic constructions.  In particular we have not fully
reproduced the asserted simultaneous stable-model and toric-uniformization
construction of Proposition 5.2, with its identification of the chosen
annular coordinate on every inverse component.  The opened Abel--Jacobi
skeleton comparison of \citet[Proposition 6.1 and Corollary 6.6]{BakerRabinoff2015}
explains the slope calculation but is not by itself the whole analytic
construction.  This is the remaining verification obligation at which we
stop promoting (D.7.1) to a theorem.  All later uses of (D.7.1) below are
explicitly conditional.  The separate attachment
\emph{\'Etale covers with a prescribed exterior sheet} is not cited by the
main section manuscript and is not used here; its theorem is claimed,
steps not checked by us.  Nor do we use Corollary 8.4, which depends on an
unattached Selmer-converse manuscript.

\statement{D.8}{Proved}{The completion assertion implies the proper local section theorem}
Assume (D.7.1) for a fixed proper hyperbolic curve $X/k$, with
$k/\mathbb Q_p$ finite.  Then $X(k)\to\operatorname{Sec}(X/k)$ is bijective.
Thus the group-theoretic and valuative passage in Section 7 of the pack does
not require an additional affine or birational-lifting conjecture.
\begin{proof}
Put $F=k(X)$, $F^{\mathrm{et}}$ as above, and prolong a rank-one valuation $w$
extending the $p$-adic valuation to $F^{\mathrm{sep}}$.
Restriction gives
\[
 D_{\bar w}(F^{\mathrm{sep}}/F)
   \longrightarrow D_{\widetilde w}(F^{\mathrm{et}}/F).
 \tag{D.8.1}
\]
It is surjective: lift an automorphism of $F^{\mathrm{et}}/F$ to
$F^{\mathrm{sep}}/F$, then correct the lift using the transitive action of
$\operatorname{Gal}(F^{\mathrm{sep}}/F^{\mathrm{et}})$ on prolongations.
The infinite-extension assertion follows from the finite assertions by
compactness.  It is injective because an element of its kernel is an isometry
fixing $F^{\mathrm{et}}$, hence fixes its completion, which contains the
specified $F^{\mathrm{sep}}$ by (D.7.1).  Normalization $w(p)=1$ ensures that
stabilizing the valuation ring preserves the absolute value itself.  The two
decomposition groups are closed, so (D.8.1) is an isomorphism of profinite groups.

No such rank-one decomposition group in $\Pi_X$ can contain the image of a
section.  Otherwise (D.8.1) lifts it to a birational section
$\widetilde s:G_k\to G_F$.  By
\citet[Proposition 2.4(2), Definition 1.4, and Observation 1.3]{Koenigsmann2005},
the entire image $\Sigma$ is also contained in a decomposition group of
$\operatorname{ord}_a$ for some $a\in X(k)$.
In the fixed field $M=(F^{\mathrm{sep}})^\Sigma$ the restrictions $v_1,v_2$
of these two selected prolongations are nontrivial rank-one henselian
valuations.  Henselianity follows because $G_M=\Sigma$ stabilizes the
prolongation, while transitivity gives uniqueness.  Algebraic extension
changes the value group by a torsion quotient and hence preserves its rank.
They satisfy $v_1(p)=1$ and $v_2|_{k^\times}=0$.

Here is the explicit contradiction, including the dyadic case.  Choose a
nonsquare $d\in k^\times$.  It is still a nonsquare in $M$, because
$\Sigma\to G_k$ is onto.  Let $t\in F$ be a uniformizer at $a$, choose an
integer $N>v_1(t)$, and set $z=p^N/t$.  Then $v_1(z)>0$ and $v_2(z)<0$.
For positive integers $r$ put
\[
 q_r=d+(1-d)\frac{z^r}{1+z^r}.
\]
The denominator is nonzero.  Direct valuation computations give
\[
 \begin{split}
 v_1(q_r/d-1)&=v_1((1-d)/d)+r v_1(z)\longrightarrow+\infty,\\
 v_2(q_r-1)&=v_2(d-1)-r v_2(z)\longrightarrow+\infty.
 \end{split}
\]
For large $r$ these values exceed $2v_1(2)$ and $2v_2(2)$, respectively.
Hensel's lemma at $1$ gives square roots of $q_r/d$ and $q_r$ in $M$;
the quotient makes $d$ a square.  The contradiction uses neither discrete
value groups nor $p\ne2$.

Now let $s:G_k\to\Pi_X$ be arbitrary.
\citet[Theorem 26 in the 2011 preprint]{PopStix2017} places its image in the
decomposition group of a valuation of $F$ extending the $p$-adic valuation.
The valuation transcendence inequality gives rational rank at most two,
hence ordered rank at most two.  Rank one has just been excluded.
For rank two, let $u$ be the rank-one coarsening.  Its prolongation is preserved
by the section: the relevant convex subgroup of the value group has a unique
extension under an algebraic field extension, because the quotient of value
groups is torsion.  If $u$ is nontrivial on $k$, the preceding rank-one
argument excludes it.  Otherwise properness gives a closed centre $a\in X$,
and a valuation dominating the discrete valuation ring $\mathcal O_{X,a}$
and trivial on $k$ is equivalent to $\operatorname{ord}_a$.
Its residue field is $k(a)$, and the remaining refinement is the unique
$p$-adic valuation of this finite extension.  The image on constants of its
horizontal decomposition group is $G_{k(a)}$.  Since $s$ projects onto $G_k$,
we obtain $k(a)=k$.  Horizontal inertia is trivial in the full \'etale tower
of the proper smooth curve, so this decomposition group is precisely the
point section at $a$.  This is also the type-$2h$ calculation in
\citet[Section 6.1]{PopStix2017}.  Surjectivity follows.

For completeness, injectivity can be checked without (D.7.1).  If $a,b\in X(k)$
give the same section, use $a$ for the Abel--Jacobi embedding into $J$.
The multiplication-by-$m$ torsor is finite \'etale for every $m\geq1$ in
characteristic zero.  Equality of section classes implies
$[b-a]\in mJ(k)$ for every $m$.  An open subgroup of the compact analytic
group $J(k)$ is $\mathbb Z_p^{g[k:\mathbb Q_p]}$.  If it has index $q$,
then $\bigcap_m mJ(k)\subseteq\bigcap_r p^r\mathbb Z_p^{g[k:\mathbb Q_p]}=0$,
using $m=qp^r$.  Thus $[b-a]=0$.  A nontrivial equivalence of the two distinct
degree-one divisors would give a degree-one map $X\to\mathbb P^1$;
this is impossible for positive genus.  Hence $a=b$.
\end{proof}

\statement{D.9}{Proved}{The exact conditional global descent statement}
Assume the proper local section theorem in D.7 over every finite completion
of a number field $K$.  For a proper hyperbolic curve $C/K$, define
\[
 C(\mathbb A_K)_\bullet=
 \prod_{v\nmid\infty} C(K_v)\times
 \prod_{v\text{ real}}\pi_0(C(K_v)).
\]
Complex places are omitted.  Let $C(\mathbb A_K)_\bullet^{\mathrm{f\text{-}cov}}$
be the locus surviving every torsor under a finite $K$-group scheme, with
survival meaning liftability to some global twist.  Then
\[
 \operatorname{loc}(\operatorname{Sec}(C/K))
   = C(\mathbb A_K)_\bullet^{\mathrm{f\text{-}cov}},
 \qquad
 C(K)\xrightarrow{\sim}\operatorname{Sec}(C/K)
 \ \Longleftrightarrow\
 C(\mathbb A_K)_\bullet^{\mathrm{f\text{-}cov}}=C(K).
 \tag{D.9.1}
\]
This is the precise conditional content of Proposition 8.1 and Theorem 8.2
of \citet{OpenAI019}.  The general equality on the right remains an additional
number-theoretic assertion, not a consequence of local point-theoreticity.
\begin{proof}
The local theorem gives a unique point at each finite place; the real section
theorem gives a connected component at each real place
\citep[Introduction, Theorem]{BrescianiVistoli2020}.  Thus localization is
well-defined.  Apply \citet[Theorem 2.1 and Remark 2.2(1), preprint
numbering]{HarariStix2012} with its normal subgroup $U=1$ and omitted set of
places $S=\varnothing$.  Its equivalence between condition (i) and condition
(iii) is exactly the equivalence between survival of all finite torsors and
interpolation of the local point sections by a global section, modulo the
full geometric group.  Evaluations are constant on real components and
complex places impose no condition.  We use the corrected convention of
\citet[Section 1]{Stoll2017Errata}; disconnected torsors are included.
This proves the first equality.

If the descent locus is $C(K)$, the local points of a section are the
localizations of one $a\in C(K)$.  Its localization has finite support
$\{a\}$, so \citet[Corollary 5]{Stix2015} makes the section geometric.
Properness excludes cuspidal sections.  Injectivity is D.5, or follows by
restriction to a single finite completion.  Conversely bijectivity of the
point-to-section map and the first equality immediately give the descent
equality.  These arguments also cover empty sets.
\end{proof}

\statement{D.10}{Proved}{What the modular-curve consequence says about a moduli section}
Assume the local theorem of D.7.  Let $N$ satisfy $g(X_1(N))\geq2$, and let
$s:G_{\mathbb Q}\to\pi_1(Y_1(N))$ be a section.  Its composite with
$\pi_1(Y_1(N))\to\pi_1(X_1(N))$ is geometrically conjugate to the section of
a unique rational cusp of $X_1(N)$.  Consequently the same conclusion applies
to a section of the arithmetic $N_{1,1}$ whose image is contained in the
open subgroup belonging to $Y_1(N)$, under the faithful stack identification
of Part E.  It does not assert that $s(G_{\mathbb Q})$ is contained in a
cuspidal decomposition subgroup of $\pi_1(Y_1(N))$.
\begin{proof}
The forgetful map from the fine modular curve $Y_1(N)$ to
$\mathcal M_{1,1,\mathbb Q}$ is finite \'etale in characteristic zero, so its
fundamental group is an open subgroup, up to conjugacy.  Containment of the
image of a section makes the same homomorphism a section of this open
subgroup's projection onto $G_{\mathbb Q}$.  Here $g(X_1(N))\geq2$ ensures
that $N$ is outside the small levels with stabilizers on $Y_1(N)$.
The compactification homomorphism is over $G_{\mathbb Q}$ and hence its
composite is a section of the proper curve.

\citet[Corollary 8.8]{Stoll2007}, with \citet{Stoll2017Errata}, proves
$X_1(N)(\mathbb A_{\mathbb Q})_\bullet^{\mathrm{f\text{-}ab}}=X_1(N)(\mathbb Q)$
when the genus is positive.  Since rational points are contained in the
finite-cover locus and that locus is contained in the finite-abelian locus,
D.9 applies.  This verifies the implication used in
\citet[Corollary 8.5]{OpenAI019} without its unrelated Corollary 8.4.
Finally \citet[Introduction, Theorem (7), p.~35]{Mazur1977} states directly
that all rational points of $X_1(N)$ are cusps whenever its genus is positive.
The equivalent torsion-order assertion is Theorem (7'); the full torsion
classification is Theorem (8).  Thus the proper section comes from a cusp,
and D.5 gives its uniqueness.

The kernel of compactification is generated normally by cusp inertia.
Identifying the image of $s$ after quotienting by this kernel does not identify
$s$ before quotienting.  D.7 is a theorem claimed for proper curves and
provides no such lifting assertion for $Y_1(N)$.  Nor does the conclusion
imply nonexistence of these splittings: a nonzero rational tangent vector at
any rational cusp gives a cuspidal section of $Y_1(N)$, exactly as in D.6.
\end{proof}

\subsection{Relative and absolute open homomorphisms}

\statement{D.11}{Proved}{An absolute open homomorphism preserves the geometric kernel}
Let $U/K$ and $V/K'$ be hyperbolic curves over number fields and let
$h:\Pi_U\to\Pi_V$ be a continuous homomorphism with open image, with no
compatibility assumed.  Then $h(\Delta_U)\subseteq\Delta_V$, and there is a
unique continuous homomorphism with open image
$\alpha:G_K\to G_{K'}$ such that $p_Vh=\alpha p_U$.
This conclusion does not depend on the 2026 open-Hom manuscript.
\begin{proof}
Set $J=p_V(h(\Pi_U))$.  The projection $p_V$ is an open map of profinite
groups, so $J$ is open in $G_{K'}$, and $J=G_L$ for a finite extension
$L/K'$ in the chosen algebraic closure.  The subgroup
$A=p_V(h(\Delta_U))$ is compact, hence closed, is topologically finitely
generated because $\Delta_U$ is, and is normal in $J$ because $\Delta_U$
is normal in $\Pi_U$.  By \citet[Theorem 2.4]{Mochizuki2006}, a number
field's absolute Galois group has no nontrivial topologically finitely
generated closed normal subgroup.  Thus $A=1$.  The universal property of
the quotient gives a unique continuous $\alpha$, and its image is $J$.
\end{proof}
The same argument is the mechanism behind
\citet[Proposition 3.3]{Mochizuki2006} for isomorphisms.  Merely calling a
geometric subgroup ``characteristic'' would not justify its preservation
by an arbitrary homomorphism; the proof above supplies the missing step.
Openness matters because it makes $J$ an absolute Galois group of a number
field, to which the theorem applies.

\statement{D.12}{Proved}{The conditional absolute Hom theorem and its known cases}
Assume \citet[Theorem 1.1]{OpenAI031} for absolute Galois groups of number
fields.  Then every $h$ as in D.11 is induced, up to target geometric inner
conjugacy after fixing the base-field map, by a unique field embedding
$j:K'\hookrightarrow K$ and a unique dominant $K$-morphism
\[
 f:U\longrightarrow V\otimes_{K',j}K.
\]
More invariantly there is a bijection
\[
 \coprod_{j:K'\hookrightarrow K}
 \operatorname{Hom}^{\mathrm{dom}}_K(U,V\otimes_{K',j}K)
 \ \simeq\
 \operatorname{Hom}^{\mathrm{open}}(\Pi_U,\Pi_V)\big/\operatorname{Inn}(\Pi_V).
 \tag{D.12.1}
\]
Without the 2026 claim the same conclusion is known whenever the induced
$\alpha$ of D.11 is injective, or is cyclotomic-compatible, and in particular
for isomorphisms $h$.  It is also known when the source field is $\mathbb Q$.
\begin{proof}
Apply the assumed theorem to $\alpha$.  It gives a unique compatible embedding
$\bar j:\overline {K'}\to\overline K$ with
$\alpha(\sigma)=\bar j^{-1}\sigma\bar j$ and $j=\bar j|_{K'}$ taking values
in $K$.  Between algebraic closures of number fields such an embedding is
onto: its image is algebraically closed and contains $\mathbb Q$, so it
contains every element algebraic over $\mathbb Q$.  Consequently $\alpha$
is injective as well as having open image.

Base change identifies the arithmetic group of $V\otimes_{K',j}K$ with
\[
 P=\Pi_V\mathbin{\times}_{G_{K'},\alpha}G_K.
\]
The map $\widetilde h(x)=(h(x),p_U(x))$ takes values in $P$ by D.11 and is
over $G_K$.  The first projection identifies $P$ with the open subgroup
$p_V^{-1}(\alpha(G_K))$ of $\Pi_V$, since $\alpha$ is injective.
Thus $\widetilde h$ has open image in $P$.  D.4 supplies the unique dominant
$f$ with the prescribed geometric conjugacy class.

We check the larger ambiguity in (D.12.1).  Conjugating $h$ by $e\in\Pi_V$
conjugates $\alpha$ by $p_V(e)\in G_{K'}$; the associated change of $\bar j$
does not alter its restriction to $K'$.  Hence $j$ is intrinsic to the absolute
outer class.  Once $\alpha$ is fixed, an $e$ conjugating one representative
to another over the same $\alpha$ has $p_V(e)$ centralizing the open subgroup
$\alpha(G_K)$.  This centralizer is trivial by
\citet[Corollary 2.2]{Mochizuki2006}.  The remaining conjugation is therefore
geometric and D.4 gives exactly the stated uniqueness of $f$.  Conversely a
pair $(j,f)$ gives such an open homomorphism: the relative map has open image
and the base-change subgroup is open in $\Pi_V$.

For the unconditional variants, an injective $\alpha$ is an isomorphism onto
an open subgroup $G_L$; the Neukirch--Uchida isomorphism theorem gives the
compatible field embedding.  One precise solvably-closed formulation is
\citet[Theorem, p.~359]{Uchida1979}, also reproduced as
\citet[Theorem 3.1]{Mochizuki2006}.  For cyclotomic compatibility use
\citet[Theorem 3.4]{Hoshi2025}.  If $h$ is an isomorphism, apply D.11 also
to $h^{-1}$, so that $\alpha$ is an isomorphism.  If $K=\mathbb Q$, use
\citet[Theorem 1]{Uchida1981}.  The rest of the proof is unchanged.
\end{proof}
For the original relative formulation in \S34, scans 483--484, the base
fields lie in one $\overline{\mathbb Q}$ and the maps are already over
$G_{\mathbb Q}$.  Its use of the Hom theorem is unconditional.  The additional
absolute assertion in D.12 is the part depending on the new claim; D.11 is
not an unresolved geometric-kernel gap.

\statement{D.13}{Open}{The 2026 unrestricted Galois open-Hom claim}
Let $E_i/F_i$ be solvably closed Galois extensions of number fields, meaning
that $E_i$ have no nontrivial abelian algebraic extensions.  The assertion
that every continuous homomorphism with open image
\[
 \alpha:\operatorname{Gal}(E_1/F_1)\longrightarrow
                \operatorname{Gal}(E_2/F_2)
\]
is induced by a unique compatible field embedding $E_2\hookrightarrow E_1$
is \citet[Theorem 1.1]{OpenAI031}, the conjecture on p.~595 of
\citet{Uchida1981}.  We record it as claimed, unrefereed, with its full
unconditional validity not completely verified by us.  D.12 is a proved
conditional consequence, not an assertion that the conjecture has thereby
been settled.

\paragraph{Proof audit of the open-Hom manuscript.}
We read Sections 1--6, including the full proofs of the new lemmas rather
than only the outline, and checked the following dependencies.
The reduction to $P=\operatorname{Gal}(\mathbb Q^{\mathrm{sol}}/\mathbb Q)$
is valid: a solvably closed algebraic field contains the roots of unity and
is closed under radicals, hence contains $\mathbb Q^{\mathrm{sol}}$.
For a number field $M\subset\mathbb Q^{\mathrm{sol}}$ this is also its maximal
prosolvable extension; a solvable normal closure of $M/\mathbb Q$ followed by
a compositum of conjugates proves the assertion.
We opened the precise local input in \citet[Section 1, pp.~595--597 and
Proposition 1]{Uchida1981}: for a fixed odd prime $\ell$, nontrivial
pro-$\ell$ inertia forces the nonabelian local case and hence an injective
map on that inertia into a unique target prime; all but finitely many target
primes occur.  The excluded torsion case is dyadic.  This result gives
cofinite coverage for each fixed $\ell$, not a bound uniform in $\ell$.
The claimed uniform bound later in the pack is therefore a real new step.
We also opened the uniqueness statement, Proposition 2 of Uchida, and
\citet[Theorem 3.4]{Hoshi2025}; the latter proves that full cyclotomic
compatibility is equivalent to geometricity.

Lemma 3.1 supplies the additive-character detection that drives the proof.
We followed its interpolation determinant and its separate translate/degree
zero estimate.  The relevant strict inequality is
\[
 n+n/r>n+1+n/m\qquad(1\leq r<n,\ m>n(n-1)).
\]
It compares the nonarchimedean determinant decay with the archimedean height
bound; no asserted Leopoldt rank is substituted.  In Proposition 3.2 the
logarithmic image of global units is the actual image, and generators of
powers of distinct prime ideals give disjoint divisor supports.  Thus the
claimed independence argument is compatible with class-field theory without
assuming Leopoldt's conjecture.  We checked the distinct roles of the
absolute-degree-one version, the Frobenius-average version, and the
nonzero-corestriction version.  Confusing them would invalidate the later
prime comparisons.

In Section 4 the map is first restricted to a normal open subgroup and then
extended to the normalizer of its kernel.  The regular-module enlargement
in Lemma 4.2 increases subgroup orders while keeping the bound on element
orders fixed.  Proposition 4.3 compares characters via a cyclic subgroup
and its proper subgroups.  We checked that its semisimple projection is
$U$-equivariant for the normal open subgroup $U$; the argument does not
require the larger intermediate subgroups to be normal in the whole
normalizer.  We followed the explicit source and target prime choices,
the residue-degree $f$, and the separation into those powers of Frobenius
that lie in $U$ and those that lie in the other kernel conjugates.
The result is surjectivity and logarithmic cyclotomic compatibility, not yet
full cyclotomic compatibility.

Section 5 passes to $\mathbb Q(i)$ and uses a sign character to show that
almost every split target prime has only finitely many possible source
partners, as the auxiliary prime varies.  Lemma 5.2 is an $S$-unit Kummer
statement; the finite class-group contribution vanishes in the relevant
inverse limit, so it is not a mistaken direct-sum assertion about all primes.
In Proposition 5.3 a finite cyclotomic twist is killed in a cyclic extension.
The extension degree can grow with $\ell$, but the multiplicity of the one
specified character in the $S$-unit representation is bounded by the number
of $S$-orbits.  The resulting bound $2[K:\mathbb Q]+2$ is independent of
$\ell$.  We checked why this independence is needed before applying the
simultaneous Kummer/Chebotarev separation of Lemma 5.4; the fixed-$\ell$
cofiniteness from Uchida cannot replace it.

Finally Section 6 compares character pullbacks at rational primes split
completely in an auxiliary finite field.  An infinite subset with one fixed
conjugate gives equality of the corresponding additive characters by
Proposition 3.2.  Faithfulness of the signed regular representation yields
equality modulo each finite quotient.  Compactness of the sets of possible
conjugating elements, with the finite-intersection property obtained by
passing to composita, supplies a single conjugation.  Extending equality
from a normal open subgroup uses the trivial centralizer of an open subgroup
of $P$.  The resulting full cyclotomic compatibility is the hypothesis of
Hoshi's theorem, which finishes the claimed proof.

These checks identify a coherent chain and expose where an apparently weaker
replacement would fail.  They do not constitute an independently reproduced
proof of every technical generality of the zero estimate in Lemma 3.1, the
character-separation argument of Proposition 4.3, and the uniform multiplicity
bound in Proposition 5.3.  We therefore do not treat those three new inputs
as established merely because their conclusions fit the reduction.  No
verified counterexample or contradictory hypothesis was found.  The attached
eight-lesson exposition is not evidence of independent verification; we do
not use it as an authority.  No Lean proof is supplied for this pack.

\subsection{Centres, restriction to covers, and the faithfulness questions}

\statement{D.14}{Known}{The centres and centralizers needed in \S35}
For a hyperbolic type $(g,n)$ in characteristic zero, the geometric profinite
surface group $\Delta_{g,n}$ is torsion-free and slim: every open subgroup has
trivial centralizer in $\Delta_{g,n}$.  In particular it and all its open
subgroups have trivial centre.  The absolute Galois group of a number field
is also slim, and hence has trivial centre and no nontrivial finite normal
subgroup.  These settle the two centre questions in \S35, folio 251 verso,
scan 496.  The second assertion does not assert that a number field's
absolute Galois group is torsion-free.
\begin{proof}
For $n>0$ the geometric group is free profinite of rank $2g+n-1\geq2$;
for $n=0$ it is the profinite completion of the closed oriented surface group
of genus $g\geq2$.
\citet[Corollaries (1.3.3)--(1.3.4)]{Nakamura1994} give centre-freeness
and the finite projective resolutions obtained from goodness of these groups.
Their cohomological dimensions at every prime are finite (one in the free
case, two in the closed surface case), so they have no elements of prime
order and are torsion-free.  Every open subgroup is again a hyperbolic
surface group of finite index, by the correspondence with finite covers.
Its Euler characteristic is its index times the original negative Euler
characteristic; thus the same centre assertion applies.

To derive slimness directly, suppose $c$ centralizes an open subgroup $H$.
Its normal core $H_0$ is open; some positive power $c^m$ belongs to $H_0$.
That power lies in the centre of $H_0$, hence is $1$.  Torsion-freeness gives
$c=1$.  This agrees with \citet[Lemma 3.2]{Mochizuki2006}.
For number fields use \citet[Corollary 2.2]{Mochizuki2006}.
If a finite subgroup $F$ is normal in a slim profinite group $G$, the kernel
of the conjugation map $G\to\operatorname{Aut}(F)$ is open and centralized
by $F$, forcing $F=1$.  Complex conjugation, when a real embedding exists,
is an element of order two of $G_K$; its subgroup is not normal.
\end{proof}

\statement{D.15}{Proved}{Rigidity on open subgroups and the two centralizer formulas}
If $G$ is slim and $u\in\operatorname{Aut}(G)$ is the identity on an open
subgroup $H$, then $u=1$.  If $u(H)=H$ and $u|_H$ is conjugation by an
element of $H$, then $u$ is that inner automorphism on $G$.
For a free basis element $x$ of $\widehat F_r$, or a standard nonseparating
simple generator $x$ of the profinite closed surface group of genus at least
two, and for every integer $m\geq1$,
\[
 C_G(x^m)=C_G(\overline{\langle x^m\rangle})
    =\overline{\langle x\rangle}.
 \tag{D.15.1}
\]
Thus Lemma 5 and equations (2), (4) in \S35, scans 498--501,
folios 252 verso--254, are valid with their stated hyperbolic hypotheses.
\begin{proof}
Let $H_0$ be the normal core of $H$ in $G$.  For $g\in G$ and $a\in H_0$
both $a$ and $gag^{-1}$ lie in $H$, so
\[
 u(g)a u(g)^{-1}=u(gag^{-1})=gag^{-1}.
\]
Consequently $g^{-1}u(g)$ centralizes the open subgroup $H_0$, hence is $1$.
This proves the first assertion.  Compose with the inverse of the specified
inner automorphism to obtain the second.

For the free basis case (D.15.1) is
\citet[Lemma (2.1.2)]{Nakamura1994}.
For the surface case, a standard generator is an algebraically simple element
in the sense of \citet[Definition 1.3]{BoggiZalesskii2019BC}.
Their Proposition 3.6 gives the stronger identity
$N_G(\overline{\langle x^m\rangle})=\overline{\langle x\rangle}$;
the centralizer contains the right-hand side and is contained in this
normalizer.  Both formulas follow.  Centralizing one element is equivalent
to centralizing the closed subgroup it topologically generates, so the
scan's alternatives $L_i^m$ and $\ell_i^m$ are mathematically equivalent in
this use.  The subgroup $\overline{\langle x\rangle}\simeq\widehat{\mathbb Z}$
has injective multiplication by every positive integer, which also completes
the power argument written across scans 500--501.
\end{proof}

\statement{D.16}{Known}{Faithful outer actions for every hyperbolic curve}
For every hyperbolic curve $U/K$ over a number field or a $p$-adic local
field, the full profinite outer representation
\[
 \rho_U:G_K\longrightarrow\operatorname{Out}(\Delta_U)
\]
is injective.  This includes affine and proper curves, type $(0,3)$,
type $(1,1)$, and all proper types $(g,0)$ with $g\geq2$.
It proves the expectation in \S35, scan 497, rather than only the weaker
assertion that a sufficiently small open subscheme is faithful.
\begin{proof}
The precise general result is
\citet[Corollary 6.3(i), with one configuration point]{HoshiMochizuki2011}.
Its full-prime-set hypothesis is essential.  The source explicitly identifies
the earlier affine assertion as
\citet[Theorem 2.1]{Matsumoto1996}; the 2011 result supplies the proper case
as well.

We also give the short Belyi argument for the particular theorem proved in
\S35, scans 491--495.  Suppose $\sigma\in G_{\mathbb Q}$ acts trivially in
$\operatorname{Out}(\widehat F_2)$ for the tripod.  Then it fixes the
isomorphism class of every connected finite \'etale cover of the tripod,
since these correspond to conjugacy classes of open subgroups of
$\widehat F_2$.  For any $j\in\overline{\mathbb Q}$ choose an elliptic curve
$E_j$ with that invariant.  By \citet[Section 4, Theorem 4]{Belyi1980} it
has a finite map to $\mathbb P^1$ branched only over $0,1,\infty$.
Deleting the inverse images of the three points gives a cover whose
$\sigma$-conjugate is isomorphic to it.  An isomorphism extends uniquely to
smooth proper completions, so $E_j\simeq E_j^\sigma$ and
$j=\sigma(j)$.  Since every algebraic $j$ was allowed, $\sigma=1$.
The unnumbered corollary on p.~256 of Belyi's English paper is exactly this
faithfulness consequence.  This argument proves injectivity, and has no
surjectivity conclusion about a separately defined normalizer or GT group.
\end{proof}
For fixed $\ell$, the analogous pro-$\ell$ assertion over a number field
is not a consequence of this theorem.  In fact for a curve with good
reduction, and a good finite place of residue characteristic different from
$\ell$, the geometrically pro-$\ell$ outer representation is unramified
there.  Nontrivial inertia of $G_K$ therefore lies in its kernel.
The full profinite statement cannot be replaced by a single-prime statement.

\statement{D.17}{Known}{Faithfulness even modulo geometric Teichm\"uller monodromy}
Let $2g-2+n>0$ and let $\mathcal M_{g,n}/\mathbb Q$ denote the stack with
ordered marked points.  Write
\[
 \Psi:\pi_1(\mathcal M_{g,n})\longrightarrow
                \operatorname{Out}(\Delta_{g,n}),\qquad
 N^{\mathrm{arith}}_{g,n}=\operatorname{im}(\Psi),\qquad
 \Sigma_{g,n}=\operatorname{im}(\pi_1(\mathcal M_{g,n,\overline{\mathbb Q}})).
\]
Then the canonical surjection
\[
 G_{\mathbb Q}\longrightarrow
       N^{\mathrm{arith}}_{g,n}/\Sigma_{g,n}
 \tag{D.17.1}
\]
is an isomorphism.  The same assertion holds with unordered marked points
and the corresponding unordered geometric monodromy group.
Here $\Sigma_{g,n}$ is the procongruence image of the orientation-preserving
mapping class group, not an assumed copy of its entire profinite completion.
This answers affirmatively the precise arithmetic-image question of \S35,
scans 502--505, including the proper types and $(1,1)$ singled out on scan 504
as possibly ``tout \`a fait faux''.
\begin{proof}
The geometric image is closed and normal because it is the continuous image
of the geometric kernel in the arithmetic stack extension.  Hence the
surjection (D.17.1) is well-defined without a rational base point on the stack.
Choose a finite extension $L/\mathbb Q$ and an $L$-point of
$\mathcal M_{g,n}$, equivalently a curve of the given type with ordered
$L$-rational marked points.  Such an $L$ exists by taking a closed point of
a nonempty finite type atlas.  We do not assume the existence of a
$\mathbb Q$-point for an arbitrary prescribed $(g,n)$.
The point gives a section over $G_L$ whose universal outer representation is
$\rho_{U/L}$.  By \citet[Corollary 6.4]{HoshiMochizuki2011},
\[
 \rho_{U/L}^{-1}(\Sigma_{g,n})=1.
\]
Thus the kernel $A$ of (D.17.1) intersects $G_L$ trivially.  It is finite,
normal in $G_{\mathbb Q}$, and therefore trivial by D.14.  This proves the
ordered assertion.

For unordered marks the ordered geometric image is a normal subgroup of
finite index in the unordered image.  Over $L$ the inverse image under
$\rho_{U/L}$ of the latter group is finite, because its map to the finite
quotient has trivial kernel by the ordered assertion of Corollary 6.4.
It is also normal in $G_L$, since geometric monodromy is normal in arithmetic
monodromy.  Slimness forces it to be trivial.  The same finite-normal-kernel
argument over $\mathbb Q$ completes the proof.
\end{proof}
The distinction from the congruence subgroup problem is exact.
\citet[Corollary 6.5]{HoshiMochizuki2011} says that the full universal
monodromy $\Psi$ is injective if and only if its geometric restriction is
injective.  Thus identifying $\Sigma_{g,n}$ with the profinite mapping class
group still requires the relevant congruence theorem.  The quotient
faithfulness (D.17.1) does not.  Equally, replacing $N^{\mathrm{arith}}$ by
a full normalizer is a different problem, retained separately in Part E.
Quotienting by orientation-reversing mapping classes would also change the
assertion; D.21 below gives the obstruction.

\statement{D.18}{Proved}{The direction of faithfulness under finite covers}
Let $k$ have characteristic zero, let $V\to U$ be a finite \'etale morphism
of geometrically connected hyperbolic curves over $k$, and assume that
$\rho_U:G_k\to\operatorname{Out}(\Delta_U)$ is faithful.  Then $\rho_V$
is faithful.  No rational point and no replacement of $k$ by a finite
extension are necessary.  This proves the finite-cover assertion in \S35,
folios 253--253 verso, scans 499--500, in its stated characteristic-zero scope.
\begin{proof}
Choose compatible base points.  Then $\Pi_V\subseteq\Pi_U$ is an open
subgroup, maps onto $G_k$ by geometric connectedness, and
$\Delta_V=\Delta_U\cap\Pi_V$ is open in $\Delta_U$.
For $\sigma\in\ker\rho_V$ choose $e\in\Pi_V$ above $\sigma$.
Its conjugation on $\Delta_V$ is conjugation by some $d\in\Delta_V$.
Thus $d^{-1}e$ centralizes $\Delta_V$.  Its conjugation on $\Delta_U$ is
an automorphism that is the identity on an open subgroup, hence is the
identity by D.14--D.15.  Therefore $\rho_U(\sigma)=1$, whence $\sigma=1$.
\end{proof}
This proof also explains the direction of the comparison map: an open
geometric subgroup detects an automorphism of the larger group once the
inner automorphism on the subgroup has been removed.  In the manuscript's
longer argument a finite residual kernel had to be eliminated; it is not
enough merely to observe that this kernel acts through a finite group.
Filling punctures is a surjection of geometric groups, so it is not the
same argument; the proper faithfulness needed after filling is supplied by
D.16 and D.17.

\statement{D.19}{Proved}{The tripod centralizer is the symmetric group}
For $U=\mathbb P^1_K\setminus\{0,1,\infty\}$ over a sub-$p$-adic field $K$,
\[
 C_{\operatorname{Out}(\Delta_U)}(\rho_U(G_K))
       =\operatorname{Aut}_K(U)\simeq\mathfrak S_3.
\]
In particular the automorphism assertion of \S34, scans 485--486, does not
require assuming a chosen arithmetic normalizer in addition to commuting
with the arithmetic action.
\begin{proof}
An outer automorphism $a$ commuting with $\rho_U$ induces, by D.1, an
isomorphism of the arithmetic extension over $G_K$.
D.4 and its Isom consequence identify it with a unique automorphism of $U$.
Conversely every such scheme automorphism commutes with the Galois action.
An automorphism extends to $\mathbb P^1$ and permutes $\{0,1,\infty\}$;
three distinct points determine a fractional linear transformation, and all
six permutations are defined over $\mathbb Q$.  This gives exactly
$\mathfrak S_3$.  Its injectivity on the outer group follows from the same
Isom theorem.  The alternatives $\mathfrak S_3\mid\mathcal S_3$ in scan 485
are two notations for this six-element group; the mathematics does not decide
the handwritten font.
\end{proof}

\subsection{Small types and the exact remaining qualifications}

\statement{D.20}{Refuted}{Dropping hyperbolicity from the faithfulness assertion}
The uniform assertion of D.16 becomes false if any of $(0,0)$, $(0,1)$,
$(0,2)$, $(1,0)$ is admitted.  The geometric centres and the counterexamples
are as follows:
\[
\begin{array}{c|c|c}
 (g,n)&\Delta_{g,n}&Z(\Delta_{g,n})\\\hline
 (0,0),(0,1)&1&1\\
 (0,2)&\widehat{\mathbb Z}&\widehat{\mathbb Z}\\
 (1,0)&\widehat{\mathbb Z}^{,2}&\widehat{\mathbb Z}^{,2}
\end{array}
\]
In contrast $(0,3)$ and $(1,1)$ are both hyperbolic with geometric group
$\widehat F_2$, but have different peripheral structures.
\begin{proof}
Over an algebraically closed field of characteristic zero, comparison with
the sphere, plane, punctured plane, and torus gives the displayed groups.
Thus the first two outer actions over $\mathbb Q$ are trivial.  For
$\mathbb G_m/\mathbb Q$ the outer action is the cyclotomic character, whose
abelian image cannot be faithful on $G_{\mathbb Q}$.  Explicitly the splitting
field of $x^3-2$ has Galois group $\mathfrak S_3$: Eisenstein's criterion at
$2$ proves irreducibility and the discriminant is $-108$, a nonsquare.
Hence $G_{\mathbb Q}$ has nontrivial commutator subgroup, killed by its
cyclotomic character.

For $(1,0)$ take $E:y^2=x^3-x$ over $K=\mathbb Q(i)$.
The endomorphism $[i](x,y)=(-x,iy)$ is defined over $K$.
Complex uniformization makes the integral lattice a free rank-one
$\mathbb Z[i]$-module, so the full Tate module is free rank one over
$\widehat{\mathbb Z}[i]$.  Equivalently one can choose a basis on which $[i]$
has matrix $J=\left(\begin{smallmatrix}0&-1\\1&0\end{smallmatrix}\right)$.
All Galois matrices commute with $J$.  Solving $AJ=JA$ over any commutative
ring gives $A=aI+bJ$; its invertible solutions commute with each other.
Thus the entire adelic representation has abelian image.  On the other hand
the $\mathfrak S_3$ splitting field of $x^3-2$ intersects $\mathbb Q(i)$
trivially: its only quadratic subfield is $\mathbb Q(\sqrt{-3})$.
Its compositum with $K$ is therefore an $\mathfrak S_3$-extension of $K$.
Consequently $G_K$ is nonabelian and the representation on
$\Delta_E=\widehat{\mathbb Z}^{,2}$ is not faithful.  The code verifies the
discriminant, the finite commutator calculation, and the centralizer
calculation modulo several integers, including $2$ and $3$; the ring
calculation above is the proof for all moduli.

For $(0,3)$ the peripheral generators are $x,y,(xy)^{-1}$, each primitive
in abelianization.  For $(1,1)$ the only peripheral subgroup is generated
by $[x,y]$, which has zero image in abelianization.  A bare identification
of their free profinite groups is therefore not an identification of curves
with their loop data.
\end{proof}

\statement{D.21}{Proved}{Orientation, real fields, and residue characteristics}
The subgroup in D.17 must be the geometric orientation-preserving group.
If a full extended mapping class group containing orientation reversal is
used instead, the analogous faithfulness assertion fails already at the
tripod.  The geometric centre statements remain valid over $\mathbb R$,
but the number-field Galois slimness statement does not extend to
$G_{\mathbb R}$.  None of D.4, D.8, D.16, or D.17 has a residue-characteristic
exclusion at $2$ or $3$ in its stated characteristic-zero setting.
\begin{proof}
Complex conjugation on the complex tripod is an orientation-reversing
homeomorphism preserving its punctures.  Its outer class belongs to the
image of the extended mapping class group and is nontrivial: on the
abelianization $\widehat{\mathbb Z}^{,2}(1)$ it acts by $-1$.
Therefore quotienting by that extended image kills a nontrivial element
of $G_{\mathbb Q}$.  Moreover the extended image need not be normal in
$N^{\mathrm{arith}}$: at the unordered tripod its image under the
isomorphism of D.17 is the subgroup generated by complex conjugation.
If it were normal it would be a nontrivial finite normal subgroup of
$G_{\mathbb Q}$, contradicting D.14.  One cannot silently form an ordinary
group quotient by it.  Taking its normal closure would still kill complex
conjugation.

Over $\mathbb R$ the arithmetic quotient is $C_2$, whose centre is itself.
The correct real section map for a hyperbolic curve is
$\pi_0(U(\mathbb R))\simeq\operatorname{Sec}(U/\mathbb R)$, by
\citet[Introduction, Theorem]{BrescianiVistoli2020}; it is not a bijection
from individual real points.  The outer action of $C_2$ is nevertheless
faithful for a hyperbolic real curve.  If the proper completion has positive
genus, conjugation reverses its orientation and is anti-symplectic on
$H_1$; identity action would contradict its nondegenerate intersection
pairing over $\mathbb Q_\ell$.  If the genus is zero and $n\geq3$, its action
on puncture homology is minus the permutation of the punctures.  An
involution on at least three punctures has at least two orbits, so the
reduced permutation representation has a nonzero $+1$ eigenspace.
The action on that subspace after the minus sign is nontrivial.

A smooth complex algebraic curve is orientable.  The fundamental group of
a nonorientable topological surface is therefore not the geometric group
of one of these curves.  A real curve can have a nonorientable quotient by
conjugation; its arithmetic datum is the geometric oriented group together
with the involution and its extension, not the fundamental group of the
quotient surface alone.  The orientation double cover does not identify
these different arithmetic extensions.

All the local fields in D.4 and D.8 have characteristic zero.  Horizontal
ramification along a deleted point is consequently tame even if its
ramification index is divisible by the residue characteristic; the analytic
unit-root threshold in D.7 handles that divisibility.  The proof of D.8
explicitly uses the correct dyadic Hensel bound.  The universal-monodromy
and full profinite theorems in D.16--D.17 are also characteristic-zero
statements.  In \S36 the orbifold orders $2$ and $3$ do prevent transporting
an argument through tame coarse moduli in characteristics $2,3$; this
requires retaining the stack and is not an exception to the theorems over
$\mathbb Q$ or $\mathbb Q_p$ stated here.
\end{proof}

\statement{D.22}{Proved}{Uniqueness of a Galois-invariant tripod loop structure}
Let $K$ be a number field and equip
$\Delta=\pi_1((\mathbb P^1\setminus\{0,1,\infty\})_{\overline K})$
with its usual outer $G_K$-action.  Among loop structures of type $(0,3)$,
meaning the three conjugacy classes
$\overline{\langle a\rangle},\overline{\langle b\rangle},
\overline{\langle (ab)^{-1}\rangle}$ for a profinite free basis $(a,b)$,
the usual loop structure is the unique one invariant under this action.
This is a precise instance of the first sentence of the conjecture in
\S34, scan 483, independent of its false normalizer explanation.
\begin{proof}
Let such a triple be invariant.  After a finite extension $L/K$, each
conjugacy class is individually fixed.  For any one of its topological
generators $z$, the normalizer of $I=\overline{\langle z\rangle}$ in
$\Pi_{U_L}$ projects onto $G_L$.  Its conjugation action on $z$ is
$z\mapsto z^{\lambda(h)}$ for some
$\lambda(h)\in\widehat{\mathbb Z}^{\times}$.
The true Galois action on $\Delta^{\mathrm{ab}}$ is scalar multiplication
by $\chi_K$, since all three actual punctures are rational.  The image of
$z$ in $\Delta^{\mathrm{ab}}$ is a primitive vector: this is evident for
$a,b$ and for $-a-b$.  Thus
$\lambda(h)[z]=\chi_K(p(h))[z]$ implies
$\lambda(h)=\chi_K(p(h))$.  In particular the geometric part of the
normalizer acts trivially.  D.3, applied over $L$, identifies $I$ with
one of the actual inertia subgroups: it is contained in one and is maximal
procyclic, being a free-basis factor or its automorphic image.
The three candidate conjugacy classes are distinct, as their abelianized
lines are distinct.  Hence they are exactly the three actual classes.
Finite base extension does not change the geometric inertia subgroups,
so the result holds over $K$.
\end{proof}

\statement{D.23}{Open}{What invariant loop data and admissible moduli splittings still require}
For an arbitrary affine hyperbolic $U/K$ over a number field, consider the
assertion that every profinite surface loop structure of its type on
$\Delta_U$ which is invariant under $\rho_U(G_K)$ equals the geometric
cuspidal loop structure.  The cyclotomic version is D.3 and the tripod case
is D.22.  We have not established this bare invariance assertion in the
remaining types, and do not identify it with the false assertion of D.2.
The first type not covered by our elementary argument is $(1,1)$: its
peripheral commutator has zero abelianization, so the primitive-vector
argument forcing the cyclotomic exponent does not apply.
Precisely the missing implication in that argument is that invariance of
an alternative peripheral subgroup forces the normalizer action required
by D.3 after an open restriction of $G_K$.

Independently, the proposed description of admissible partial splittings
of $N_{g,n}$ in \S34, folios 247--247 verso, scans 487--488, is not a
consequence of the Hom theorem or the proper section manuscript.
Once $N^{\mathrm{arith}}$ is identified with an arithmetic moduli-stack
fundamental group in a type where the congruence kernel vanishes, these
splittings concern points of that stack.  Classifying all of them as
curve-moduli points is a moduli-stack section problem; neither D.5 nor
D.7 supplies that classification.  In particular the empty list of
extra conditions at $(1,1)$ in scan 487 is not proved sufficient by the
2026 proper-curve result.  D.10 is the exact consequence obtained on the
specified high-genus modular-curve open subgroup.  In higher genus, even
the identification with the full arithmetic stack fundamental group must
retain the congruence-kernel qualification of D.17.

These Open labels state the exact assertions left unproved by this audit;
they are not claims that every formulation of inertia recognition or every
special section theorem is unknown in the literature.  Our searches and
checks used Nakamura's cyclotomic criterion, Mochizuki's relative Hom theorem,
the full Hoshi--Mochizuki monodromy theorem, and the two new proof packs.
An assertion left conditional or unverified has not been refuted merely
because our proof stops at an identified implication.
